# 📘 PART I — Front Matter, Foundations, Geometric Layer

`latex
\documentclass[11pt]{article}
\usepackage{amsmath, amssymb, amsthm}
\usepackage{geometry}
\usepackage{hyperref}
\geometry{margin=1in}

\title{Anima-Vectorium Mathématique v5.0 \\
\large Geometric–Thermodynamic Foundations for the Vectorium Engine}
\author{Borealis S. Hedling}
\date{\today}

\begin{document}
\maketitle

\begin{abstract}
Anima-Vectorium v5.0 establishes a unified mathematical backbone integrating
geometric curvature, holonomy transport, non-linear operator algebra,
thermodynamic free-energy dynamics, and spectral-to-geometric translation rules.
This document supersedes v4.0 and provides the complete mathematical foundation
required for executable Rust/Python/FFI implementation of the Vectorium engine.
\end{abstract}

\tableofcontents

%------------------------------------------------------------
\section{Foundations}

\subsection{State Space}
The Vectorium state space is defined as
\[
\mathcal{S} = \{ (x, M, \kappa, \delta) \mid x \in \mathbb{R}^n \},
\]
where $x$ is the tensor state, $M$ is metadata, $\kappa$ is the curvature field,
and $\delta$ is the drift field.

\subsection{Tensor Manifold}
The tensor manifold $\mathcal{M}$ is a smooth Riemannian manifold equipped with
a metric $g$. For tangent vectors $u, v \in T_x\mathcal{M}$,
\[
g_x(u, v) = u^\top G(x) v,
\]
where $G(x)$ is a positive-definite matrix field. The Euclidean case corresponds
to $G(x) = I$. Vectorium-specific curvature behavior is modeled by
\[
G(x) = I + \lambda \cdot \Phi(x),
\]
where $\Phi(x)$ is activation-induced curvature.

\subsection{Metadata Fields}
Metadata $M$ consists of symbolic labels and scalar weights that do not
contribute to curvature. Metadata is preserved under holonomy transport and
does not affect thermodynamic gradients.

%------------------------------------------------------------
\section{Geometric Layer}

\subsection{Curvature Operator}
Curvature is defined by
\[
\kappa(x) = \nabla^2 x + \Phi(x),
\]
where $\Phi(x)$ is activation-induced curvature (Section 3). The Laplacian
$\nabla^2$ is defined with respect to the metric $g$.

\subsection{Connection for Parallel Transport}
Holonomy transport requires a connection $\nabla$ compatible with $g$. Vectorium
uses the Levi-Civita connection:
\[
\nablau v = \partialu v + \Gamma(u, v),
\]
with Christoffel symbols
\[
\Gamma^k_{ij} = \frac{1}{2} G^{kl}
\left(
\partiali G{jl}
+ \partialj G{il}
- \partiall G{ij}
\right).
\]

\subsection{Holonomy Transport}
Parallel transport along a curve $\gamma(t)$ satisfies
\[
\nabla_{\dot{\gamma}(t)} v(t) = 0.
\]
Holonomy around a closed loop $\gamma$ is
\[
\mathcal{H}(x) = x + \oint_{\gamma} \kappa \cdot d\gamma.
\]

\subsection{Drift Fields}
Drift fields $\delta$ represent directional tendencies in state evolution.
Spectral drift fields are translated into geometric gradient fields (Section 5).

\subsection{Geometric Invariants}
Curvature invariants:
\[
\int \kappa(x) \, dx = \text{constant}.
\]
Holonomy invariants:
\[
\oint_{\gamma} \kappa \cdot d\gamma = \text{constant}.
\]
`

---

